\textbf{Clutter is the clearest degrading factor in this design.}
Increasing obstacles from 0 to 5 raises $L$ by $9.5\pm3.0$ steps; its
95\% interval $[6.4,12.5]$ excludes zero. Its effect exceeds the state-size
effect by 11.5 steps (95\% CI $[7.1,15.9]$) and the action-size effect by
19.7 steps (95\% CI $[15.9,23.5]$). This is evidence for a ranking within
these tested levels, not a claim that obstacle count always dominates every
other source of difficulty. Obstacles affect both routes and the supplied
distance-dependent noise.

\textbf{State count alone does not explain online difficulty.}
Increasing $n_y$ from 5 to 20 changes $L$ by $-2.0\pm2.3$ steps,
with interval $[-4.4,0.3]$. There is no clear degradation. Mean setup time
increases from 18.1 to 63.2\,ms because goal-distance maps cover more cells,
but search does not enumerate the $F^2$ joint configurations. The wider
workspace also offers more recovery space, so this comparison cannot isolate
an abstract state-count effect from all geometric changes.

\textbf{More available actions improve these tasks.}
Increasing the joint intended-action count from 36 to 676 changes $L$ by
$-10.2\pm2.2$ steps, with interval $[-12.5,-8.0]$. Diagonal commands can
combine forward progress with recovery from lateral/vertical slips, and
the heuristic's deterministic relaxation is closer to the richer action set.
These are plausible explanations, rather than separately identified causes.
The mean number of admitted root actions grows only from 5.8 to 6.4;
widening prevents the coded action count from translating directly into full
enumeration. The experiment therefore measures the net effect of additional
control choices and search cost, not a pure branching-factor penalty.
